\section{A closed stratum and compatible algebraic products}\label{sec:stratum}

We turn the free part of the local cotangent module into a closed
geometric stratum, then construct algebraic product presentations along
it.  The transition maps preserve its marked ideal at every order.
This supplies a single system in which the ambient tensor maps and the
coherent coefficients supported on finite thickenings can be compared.

For a normal complex analytic germ $(V,v)$ put
\[
 D_v=\im\bigl(T_V{}_v\longrightarrow T_{V,v}^{\mathrm{Zar}}\bigr),
 \qquad r(v)=\dim_{\C}D_v.
\]
Here the arrow evaluates a germ of a derivation at $v$, and
$T_V=\Hom(\Omega_V^1,\cO_V)$.  All germs in the next two lemmas are finitely
dominated by smooth germs of the same dimension, as supplied by
Proposition~\ref{loc:maximal-charts}.

\begin{lemma}[Evaluation and smooth factors]\label{str:evaluation}
The largest rank of a free direct summand of $\Omega_V^{[1]}{}_v$ is
$r(v)$.  There is an analytic product decomposition
\begin{equation}\label{str:analytic-product}
 (V,v)\simeq(\Delta^{r(v)},0)\times(T,t),
\end{equation}
where all vector fields on $(T,t)$ have zero evaluation at $t$.
Moreover, if $\beta$ is a reflexive one-form and a vector field $\xi$
has zero evaluation at $v$, then $\beta(\xi)(v)=0$.
\end{lemma}

\begin{proof}
We first prove the last assertion, which is stronger than the analogous
statement for ordinary differentials.  A derivation has a local holomorphic
flow: after embedding the germ in a polydisc, lift its coordinate values to
an ambient vector field and solve the holomorphic differential equation.
Preservation of the defining ideal makes the flow preserve the embedded
analytic space.  If the derivation evaluates to zero, its flow $\phi_z$
fixes $v$.

Choose a finite smooth cover $a:(Y,y)\to(V,v)$ and consider
$H(z,w)=\phi_z(a(w))$.  The map $(z,w)\mapsto(z,H(z,w))$ is finite:
it is $\id_\Delta\times a$ followed by the flow automorphism of
$\Delta\times V$.  Thus the inverse image of $\Delta\times V_{\mathrm{sing}}$
has codimension at least two.  The pullback of $\beta$ from
$V_{\mathrm{reg}}$ consequently extends to a holomorphic one-form
$\alpha$ on the smooth space $\Delta\times Y$.

We claim that $\alpha$ restricts to zero on $C=\Delta\times\{y\}$.
To check this, pull a fixed finite smooth cover $a_0:Y_0\to V$ back by
$H$, normalize a component dominating $\Delta\times Y$, and resolve it.
The resulting diagram has a proper dominant map $q$:
\[
 \begin{array}{ccc}
 \widetilde W&\xrightarrow{\ g\ }&Y_0\\
 {\scriptstyle q}\downarrow&&\downarrow{\scriptstyle a_0}\\
 \Delta\times Y&\xrightarrow{\ H\ }&V.
 \end{array}
\]
The pullback of $\beta$ extends to a holomorphic form $\beta_0$ on
$Y_0$.  On $\widetilde W$ the forms $q^*\alpha$ and $g^*\beta_0$
agree, since they agree on a dense open.  An irreducible subvariety of
$q^{-1}(C)$ dominating $C$ maps under $g$ to a point of the finite fiber
$a_0^{-1}(v)$.  Restriction to a resolution of that subvariety therefore
annihilates $q^*\alpha$.  Pullback of one-forms from the smooth curve
$C$ under a dominant map is injective in characteristic zero, proving
the claim.  Finally, the definition of the flow gives the identity
\[
 \left.\alpha\left(\frac{\partial}{\partial z}\right)\right|_{z=0}
       =a^*(\beta(\xi)).
\]
It holds first over the regular locus and then everywhere by holomorphic
extension.  Evaluation at $y$, using $\alpha|_C=0$, yields
$\beta(\xi)(v)=0$.

Suppose that $\Omega_V^{[1]}{}_v$ has a free summand with basis
$\beta_1,\ldots,\beta_s$ and dual vector fields
$\xi_1,\ldots,\xi_s$.  A linear combination of their evaluations that
vanishes pairs to zero with every $\beta_j$ by the preceding paragraph.
The identities $\beta_j(\xi_i)=\delta_{ij}$ imply independence of the
evaluations.  Thus $s\leq r(v)$.

Conversely choose vector fields with $r=r(v)$ independent evaluations and
regular functions $z_1,\ldots,z_r$ with
$\xi_i(z_j)(v)=\delta_{ij}$.  Cut a transverse germ by $z_1=\cdots=z_r=0$.
Successive flows of the chosen fields define a map from a polydisc times
this transverse germ to $V$.  To invert it, apply the inverse flows and
solve for their times by the analytic implicit function theorem; its
Jacobian in the time variables is the identity at $v$.  This argument
works in an ambient embedding and hence also on the embedded germ.
The resulting product has $r$ free cotangent summands.  The transverse
factor cannot have a vector field with nonzero evaluation, since such a
field would give an additional independent evaluation on the product.
\end{proof}

\begin{lemma}[The rank-locus dimension bound]\label{str:rank-locus}
On such a normal germ, the locus $\{v:r(v)\leq a\}$ is closed analytic
and has dimension at most $a$.
\end{lemma}

\begin{proof}
Coherence of the sheaf of derivations, followed by evaluation on finitely
many ambient coordinates, expresses the condition by minors of a matrix
of holomorphic functions.  For the dimension assertion we use the
following trace observation.  If $\pi:(Y,y)\to(V,v)$ is a finite map
from a smooth germ, with this single point over $v$, and $m$ is its degree,
then a vector field $\eta$ on $Y$ induces the derivation
\begin{equation}\label{str:trace-derivation}
 a\longmapsto\frac{1}{m}\operatorname{Tr}_{\C(Y)/\C(V)}
                    \bigl(\eta(\pi^*a)\bigr).
\end{equation}
The trace is holomorphic by integrality and normality, and the Leibniz
rule follows from linearity of trace over $\cO_V$.  At $v$, the trace of a
holomorphic function on this local branch is $m$ times its value at $y$.
Consequently the evaluation of \eqref{str:trace-derivation} is exactly
$d\pi_y(\eta_y)$.  A local branch can be isolated analytically even if
the original finite cover has several points over $v$.  We obtain
\begin{equation}\label{str:differential-image}
 d\pi_y(T_{Y,y})\subset D_v.
\end{equation}

Let $Z$ be an irreducible subvariety of the rank locus.  A component of
its reduced inverse image under a finite smooth cover dominates $Z$.
At general smooth points of this component and of $Z$, the restricted
finite map has differential rank $\dim Z$, by characteristic zero.
Equation~\eqref{str:differential-image} therefore gives $\dim Z\leq a$.
\end{proof}

Define $s(b)$ to be the free cotangent rank on a maximal chart at $b$.
It is independent of the choices in Proposition~\ref{loc:maximal-charts}; it
is also the multiplicity of the tensor unit in Theorem~\ref{loc:slsp}.

\begin{proposition}[The minimum stratum]\label{str:minimum}
Let $d=\min_{b\in B}s(b)$.  If $d=n$, every germ of $B$ is a quotient
singularity.  Otherwise the minimum locus is a nonempty closed algebraic
subset.  Each of its connected components $S_0$ is irreducible, normal
and projective of dimension $d$.  On a maximal chart over a point of
$S_0$, its reduced inverse image is, as a germ of a pair,
\begin{equation}\label{str:minimum-pair}
 (V,S_V)\simeq(\Delta^d\times T,\Delta^d\times\{t\}).
\end{equation}
At every point of $S_V$, the vector-field evaluation space is exactly
the tangent space to $S_V$.
There is a smooth separated Deligne--Mumford stack $\cS$ with coarse
space $S_0$ whose charts are these smooth factors.  Its stabilizers are
the full deck stabilizers, including those acting trivially on the
smooth factor.
\end{proposition}

\begin{proof}
Smooth factors persist under further finite local quasi-\'etale covers.
Indeed such a cover is \'etale over the smooth locus by purity, and a
cover of $\Delta^r\times T_{\mathrm{reg}}$ is pulled back from
$T_{\mathrm{reg}}$, since the polydisc is simply connected.  Normalization
extends the product across the omitted codimension-two set.  More
precisely, if $c'$ is any nearby point on a finite quasi-\'etale chart
over $b'$, its pointed local branch is dominated by a maximal cover of
$(B,b')$.  Each of its $r(c')$ smooth factors persists on that cover,
so $r(c')\leq s(b')$.  A product with $r$ smooth directions therefore
implies $s\geq r$ throughout its image.  Each pointed finite branch
dominates its base germ, so the image of a sufficiently small
neighborhood contains a neighborhood downstairs.  Along the smooth factor
through its marked point, the transverse germ is unchanged; the same
description of covers shows that its maximal-cover property is unchanged
as well.  Hence $s$ is lower semicontinuous in the analytic topology
and is constant along that factor.

At a minimum point choose the product in Lemma~\ref{str:evaluation}.
The preimage of the minimum locus lies in the evaluation-rank locus
$r\leq d$: any extra smooth factor at another point of this chart
would persist on its maximal cover.  The rank locus on the full chart
is the product of $\Delta^d$ with the transverse rank-zero locus $Z$.
Apply Lemma~\ref{str:rank-locus} to this full chart, which is finitely
dominated by a smooth germ.  It gives $d+\dim Z\leq d$, so the germ of
$Z$ at the marked point is just that point.  Conversely the whole smooth factor
has the original maximal transverse type.  This proves
\eqref{str:minimum-pair} and also proves that the minimum locus is locally
analytic, rather than merely topologically closed.

Its closedness and projectivity of $B$ make it algebraic by Chow's
theorem.  Locally it is a finite quotient of the smooth orbit by its
deck stabilizer, and hence is normal.  Distinct irreducible components
of a normal variety are disjoint, so a connected component is
irreducible.  If $d=n$, Lemma~\ref{str:evaluation} makes every maximal
chart smooth, which is exactly the quotient-singularity assertion.
The construction of the stack, with its specified stabilizers, is
included in Proposition~\ref{str:products} below.
\end{proof}

From now on fix such a component with $d<n$, and put $c=n-d>0$.
We first prove the precise algebraization statement needed for its
neighborhood.

\begin{lemma}[Algebraizing a marked product]\label{str:marked-artin}
Let $(V,v)$ be an algebraic complex germ, let $S\subset V$ be a reduced
closed subgerm, and suppose the analytic pair is
$(\Delta^d\times T,\Delta^d\times\{t\})$.  There are an algebraic
transverse germ $(T_0,t_0)$ and common \'etale neighborhoods of pairs
\begin{equation}\label{str:etale-product-pair}
 (V,S)\ \longleftarrow\ (U,S_U)\ \longrightarrow\
 (\mathbb A^d\times T_0,\mathbb A^d\times\{t_0\}).
\end{equation}
Both inverse-image ideals defining $S_U$ are equal, without truncation.
For $k\geq1$, put $A_k=\cO_{T_0,t_0}/\mathfrak m_{t_0}^k$ and let $U_k$
be defined by the $k$th power of this common ideal.  After shrinking the
displayed neighborhoods, there are compatible isomorphisms
\begin{equation}\label{str:finite-product}
 U_k\simeq S_U\times\Spec A_k
 \quad\text{over}\quad \mathbb A^d\times\Spec A_k.
\end{equation}
No equivariance of these product isomorphisms is asserted.
\end{lemma}

\begin{proof}
Choose algebraic functions $z_1,\ldots,z_d$ whose restrictions give
independent cotangents on $S$ at $v$, and define $T_0$ by their vanishing.
In the given analytic product these functions may replace the polydisc
coordinates: their Jacobian in those coordinates is invertible.  The
analytic implicit function theorem, with the transverse germ retained,
identifies the resulting transversal with $T$ and gives a formal
isomorphism of the algebraic pointed pairs in
\eqref{str:etale-product-pair}.

Here is why approximation can retain the marked ideal exactly.  Write
finite algebra presentations for the target product and generators
$a_1,\ldots,a_h$ and $b_1,\ldots,b_l$ for its marked ideal and the ideal
of $S$ in $V$, respectively.  Besides the target's defining equations,
include unknown coefficients in the finitely many equations
\[
 \varphi(a_i)=\sum_j u_{ij}b_j,
 \qquad b_j=\sum_i v_{ji}\varphi(a_i).
\]
The formal pair isomorphism supplies a solution.  Artin approximation
over the henselization of the algebraic local ring produces a solution
in an \'etale neighborhood, arbitrarily close to that formal solution;
this is the polynomial-equation form of
\cite[Tags 0CAU and 0CAV]{StacksProject}.  Denominators in local
presentations can be included by equations for their inverses.  The
displayed equations give the two ideal inclusions exactly.

Choose approximation modulo the square of the maximal ideal.  Composing
the resulting map on completions with the original formal inverse
gives an endomorphism that is the identity on residue field and
cotangent space.  It is the identity on the associated graded ring,
which is generated in degree one.  Induction on the powers of the
maximal ideal and completeness show that it is an isomorphism.  Thus
the second map in \eqref{str:etale-product-pair} induces an isomorphism
of completed local rings and is \'etale at the marked point; shrink
to its \'etale locus.  This verifies the inverse-map issue even for a
singular transversal.

The restriction $S_U\to\mathbb A^d$ is \'etale.  Both $U_k$ and
$S_U\times\Spec A_k$ are \'etale over
$\mathbb A^d\times\Spec A_k$ and have the same reduction over
$\mathbb A^d$.  Uniqueness of lifting morphisms between \'etale schemes
through a nilpotent thickening gives \eqref{str:finite-product}
\cite[Tag 025H]{StacksProject}.  The same uniqueness makes the
isomorphisms compatible as $k$ varies.  Notice that this last argument
uses the equality of the two ideals, not closeness of those ideals to
some finite order.
\end{proof}

\begin{proposition}[Compatible product presentations]\label{str:products}
The following data exist along $S_0$.
\begin{enumerate}
\item There are finitely many normal finite Galois quasi-\'etale charts
$V_\alpha\to B_\alpha$, with deck groups $\Gamma_\alpha$ and with
$B_\alpha\to B$ separated and \'etale, covering $S_0$.
They are maximal at every point over $S_0$.
\item Normalized main overlaps give an \'etale groupoid along the
reduced inverse images $S_\alpha\subset V_\alpha$.  Its quotient is
the stack $\cS$ of Proposition~\ref{str:minimum}.  Writing $I$ for the
marked ideal on the object charts, the same finite diagrams define,
for every $k\geq1$, a proper separated Deligne--Mumford stack
$\mathcal S_k$, with $\mathcal S_1=\cS$.  Its coarse space is finite
over $B$ and projective.
\item A finite family of \'etale refinements of the object charts
covers their closed strata and has the marked-product presentations of
Lemma~\ref{str:marked-artin}.  All transition maps come from the
normalized ambient overlaps and preserve $I$ exactly.  Their
restrictions are consequently compatible for all powers $I^k$.
\item The normalized source charts give ordinary pullback and descent
of coherent sheaves on any $\mathcal S_k$ to coherent sheaves on $X$
supported over $S_0$.
\end{enumerate}
These assertions are encoded by finitely many algebraic schemes,
morphisms, ideals and identities.  In particular they can be spread
simultaneously before choosing any order of Frobenius pullback.
\end{proposition}

\begin{proof}
Start with a chart maximal at one point, as in
Proposition~\ref{loc:maximal-charts}.  By
\eqref{str:minimum-pair} and Lemma~\ref{str:marked-artin}, its transverse
analytic type is constant along the part of $S_\alpha$ covered by one
common \'etale product neighborhood.  The image is open in $S_\alpha$,
because an \'etale map is open.  Every point there still has a maximal
germ, since the transverse germ and its covers have not changed.
Take the union of translates of this open neighborhood around the finite
deck orbit over the base point.  The complement in $S_\alpha$ has closed
image under the finite chart map.  Removing that image downstairs
arranges maximality at every point over the remaining stratum.
Quasi-compactness of $S_0$ gives the finite list in (1).  Applying
Lemma~\ref{str:marked-artin} at its inverse images and using
quasi-compactness gives a finite family of product refinements as in (3).

We construct the groupoids while keeping the full finite normalizations.
Put $V=\coprod_\alpha V_\alpha$.  Choose a common dense open
$B^\circ\subset B_{\mathrm{reg}}$ over which all the charts are \'etale,
and write $V^\circ=V\times_BB^\circ$.  For $a=2,3,4$, let $W_a$ be the
reduced union of the irreducible components of $V^{\times_B a}$ meeting
$(V^\circ)^{\times_{B^\circ}a}$, and let
\[
 \overline V_a=W_a^\nu
\]
be its full normalization.  It is finite over the closed subscheme
$W_a\subset V^{\times_Ba}$.  Its restriction over $B^\circ$ is the
ordinary product, which is \'etale over the smooth space $B^\circ$ and
therefore normal.  Write $\overline R=\overline V_2$, with projections
$s,t:\overline R\to V$.

The projections of these normalizations are \'etale at every point
over $S_0$.  To see this locally, pass to a common \'etale base
neighborhood for the chosen object charts.  A projection is there a
finite quasi-\'etale map followed by an \'etale map: away from the
codimension-two branch loci it is a base change of the original charts,
and finiteness preserves that codimension bound.  Each pointed target
chart over $S_0$ is maximal, so every connected local finite
quasi-\'etale factor is trivial.  This proves the asserted \'etaleness.

We will also need an arrow through every pair of object points over a
point of $S_0$.  Pass to the common strict henselian base germ at that
point and select the pointed local factors of the two finite charts.
By the algebraic--analytic comparison in
Proposition~\ref{loc:maximal-charts}, the maximality at these points
also gives maximality of the selected strict henselian factors.
They are therefore isomorphic over this germ.
The graph of such an isomorphism is a main branch of their product and
passes through the chosen pair.  Normalization commutes with this
strict henselian base change: the base change of a finite normalization
is normal under \'etale localization and has the same generic
normalization property, and the assertion passes to the filtered
strict henselian limit.  The graph branch therefore supplies a point
of $\overline R$ over the chosen pair.  This proves the assertion for
the original algebraic charts.

Let $R\subset\overline R$ be the open where both $s$ and $t$ are
\'etale.  It contains the entire inverse image of $S_0$.  Let
$I_\alpha$ be the reduced ideal of $S_\alpha$ and let $I$ be their
union on $V$.  On the full normalization define the extended ideals
\[
 J_s=\im(s^*I\longrightarrow\cO_{\overline R}),\qquad
 J_t=\im(t^*I\longrightarrow\cO_{\overline R}).
\]
Their zero sets are both the inverse image of $S_0$.  On $R$ they are
equal: the two \'etale pullbacks are reduced and have that same support.
Let $V_k\subset V$ be defined by $I^k$, and define $R_k$ to be the
closed subscheme of the full $\overline R$ defined by $J_s^k$.
Its entire support lies in $R$, so
\[
 R_k=R\times_{s,V}V_k=R\times_{t,V}V_k.
\]
In particular its two maps to $V_k$ are \'etale.  The map to
$V\times_BV$ factors through $V_k\times_BV_k$, because both extended
ideals vanish to order $k$ on $R_k$.  The resulting map
$R_k\to V_k\times_BV_k$ is finite: it is a closed quotient of the
full finite normalization, with a factorization through that closed
subscheme.  This is the finiteness used below.

Here are the normalization comparisons needed for composition.  Put
\[
 P=R\times_{t,V,s}R.
\]
It is normal, since it is \'etale over the normal object charts.
Every irreducible component of $P$ meets the open over $B^\circ$:
its image under an \'etale projection to an object component is a
nonempty open, and therefore meets the dense $V^\circ$.  Over
$B^\circ$, the scheme $P^\circ$ is an open subscheme of the ordinary
triple product $(V^\circ)^{\times_{B^\circ}3}$.  Hence each component
of $P$ maps dominantly onto a main component of $W_3$.  Its generic
map lifts to that component of $\overline V_3$, and the lift extends
over the normal component of $P$.  Indeed, each element of the
finite normalization is integral over the local target ring; its
generic pullback is integral over a normal local ring of $P$ and so
belongs to that ring.  This constructs $P\to\overline V_3$.

The three pair-forgetting maps $\overline V_3\to W_2$ lift in the
same way to maps $\rho_{01},\rho_{12},\rho_{02}:\overline V_3\to
\overline R$.  Componentwise dominance here follows on $B^\circ$:
each projection of the ordinary triple product is \'etale, and its
nonempty image is open in the corresponding ordinary pair component.
Let
\[
 O_3=\rho_{01}^{-1}(R)\cap\rho_{12}^{-1}(R).
\]
It contains the entire inverse image of $S_0$, and its two adjacent
pairs give a map $O_3\to P$.  Conversely, the adjacent pair maps of
$P\to\overline V_3$ agree with the two projections $P\to R$ on the
dense unbranched open, hence on $P$; the lift therefore lands in $O_3$.
This map and $O_3\to P$ are
inverse between $O_3$ and $P$.  Their composites are the identity
on the componentwise dense ordinary unbranched products; equality
extends because the sources are reduced and the targets separated.
Restrict further to
$O'_3=O_3\cap\rho_{02}^{-1}(R)$, which still contains the entire
closed support.  The map $\rho_{02}$ gives composition there.

On $P$ the three extended ideals from the object coordinates are
equal, by the equality $J_s=J_t$ on each arrow.  Call their common
ideal $J$.  Base change by the closed immersions therefore gives the
scheme equality
\[
 R_k\times_{V_k}R_k=P/J^k.
\]
Its support is the whole inverse image of $S_0$ in $P$, hence lies in
the domain corresponding to $O'_3$.  The three ideals agree there
also along $\rho_{02}$, so composition restricts to
$R_k\times_{V_k}R_k\to R_k$ for every $k$.

For associativity use $R\times_VR\times_VR$ and $\overline V_4$.
The same argument applies component by component: the former is
normal, and its unbranched part is open in the ordinary quadruple
product and meets every component.  Each component therefore dominates
a main component of $W_4$, while all pair-forgetting projections are
dominant onto their main pair components there.  It identifies with the open of
$\overline V_4$ where the adjacent pairs $01,12,23$ land in $R$.
Restrict to the further open where $02,13,03$ also land in $R$.
Every one of these opens contains the entire inverse image of $S_0$.
Both parenthesizations are defined on this last open and agree on
its componentwise dense unbranched product, hence agree everywhere.
Their restrictions agree on every quotient by the common $k$th
power of the stratum ideal.  Interchanging the two coordinates lifts
to the normalization $\overline R$ and gives inversion on the open
where both the original and interchanged arrow land in $R$.  The
diagonal of the ordinary unbranched pair product is open and closed;
its main branch lifts the diagonal map from the normal $V$ to
$\overline R$.  Restricting to the inverse image of $R$ gives the unit.
These domains contain the entire closed support, by the \'etaleness of
all pair points over $S_0$.  The unit and inverse identities follow on their
support-containing domains by the same dense-open comparison and
then restrict to all $k$.  Thus $R_k\rightrightarrows V_k$ is an
\'etale groupoid.  Four factors suffice because associativity is the
only law involving three composable arrows.

Since $B$ is separated, $V_k\times_BV_k$ is closed in $V_k\times V_k$.
The finiteness proved above therefore makes the diagonal of
$[V_k/R_k]$ finite.  This quotient is a separated Deligne--Mumford
stack with finite inertia.  Its local finite quotient presentation
can be read after base change to $B_\alpha$.  The diagonal
$B_\alpha\to B_\alpha\times_BB_\alpha$ is open and closed, since
$B_\alpha\to B$ is separated and \'etale.  On that component the
self-overlap is the normalization of
$V_\alpha\times_{B_\alpha}V_\alpha$, namely
$\coprod_{\gamma\in\Gamma_\alpha}V_\alpha$ with its distinct deck
labels.  The existence of a main arrow through every pair of closed
object points, proved above, makes the other object charts \'etale
refinements of this presentation: the relevant arrow projections
are \'etale and surjective on the closed support, which is the
underlying support of every $V_k$.  Consequently over $B_\alpha$
the stack is
\[
 [\,(V_\alpha)_k/\Gamma_\alpha\,].
\]
Its coarse algebra is precisely
\[
 (\cO_{V_\alpha}/I_\alpha^k)^{\Gamma_\alpha}
   =\cO_{B_\alpha}/(I_\alpha^k\cap\cO_{B_\alpha}).
\]
Here the intersection is the contracted ideal in the finite chart
algebra; exactness of finite-group invariants in characteristic zero
gives the equality.  These finite algebras glue \'etale locally.
They give a coarse space finite over the projective variety $B$, and
therefore projective.  A quotient by a finite group is proper over
its coarse space, so the stack is proper.  At $k=1$ the coarse space
is $S_0$ and the object charts are smooth, giving $\mathcal S_1=\cS$.

The construction keeps the label of every $\gamma\in\Gamma_\alpha$
even when it fixes the whole smooth factor.  Stabilizers may increase
on closed subsets of $S_\alpha$, and remain subgroups of these fixed
finite deck groups.  The ambient groupoid laws and their restrictions
therefore retain the full stabilizers on $\cS$.

For (4), write $X_\alpha\to V_\alpha$ for the normalized source chart.
Proposition~\ref{loc:maximal-charts} makes
$X_\alpha\to X\times_BB_\alpha$ a finite \'etale
$\Gamma_\alpha$-torsor.  The normalized overlaps give its ordinary
descent identifications.  The ideals
$I_\alpha^k\cO_{X_\alpha}$ descend to a coherent ideal near
$f^{-1}(S_0)$; gluing it with the unit ideal off this closed set gives
a closed subscheme $X_k\subset X$.  This is the descended ideal
quotient, with no identification with a fiber product over the coarse
thickening required.  Put $X_{\alpha,k}=X_\alpha/I_\alpha^k\cO_{X_\alpha}$.
The local diagram
\[
 \begin{array}{ccc}
 X_{\alpha,k}&\longrightarrow&(V_\alpha)_k\\
 \downarrow&&\downarrow\\
 {[X_{\alpha,k}/\Gamma_\alpha]}&\longrightarrow&
       [(V_\alpha)_k/\Gamma_\alpha]
 \end{array}
\]
is $2$-Cartesian.  Its lower left quotient is the descended scheme
$X_k\times_BB_\alpha$, because the source action is a torsor.
These diagrams glue to a map $f_k:X_k\to\mathcal S_k$.
For a coherent sheaf $M$ on $\mathcal S_k$, the ordinary sheaves
$\cO_{X_{\alpha,k}}\otimes_{\cO_{(V_\alpha)_k}}M_\alpha$ descend to
a coherent sheaf on $X_k$, and then on $X$.  Increasing $k$ does not
change this pullback, denoted $f^*M$.  The displayed Cartesian
diagram is also the reason that ordinary chart base change and the
canonical higher-direct-image coefficient maps descend.  After the
product refinements in (3), these constructions are their ordinary
\'etale pullbacks from the finite charts.

All schemes, full finite normalizations, group actions, product maps,
marked ideals and maps in the diagrams with at most four factors are
of finite type.  Their identities, ideal inclusions, and the
disjointness of each omitted closed complement from its closed support
are finitely many algebraic conditions.  Individual marked pairs were
approximated in Lemma~\ref{str:marked-artin}; their transition maps are
the normalized algebraic maps just constructed.  This proves the
asserted finite nature of the compatible data.
\end{proof}

We write $\widehat{\mathcal B}$ for the compatible formal system
$(\mathcal S_k)_{k\geq1}$.  There are two kinds of sheaves in use on this
system.  First, an ambient reflexive tensor object
$\mathcal M_\alpha$ on each normal chart, with its natural \'etale
transition isomorphisms, gives the completed system
\[
 M=(M_k)_{k\geq1},\qquad
 M_{\alpha,k}=\mathcal M_\alpha/I_\alpha^k\mathcal M_\alpha.
\]
In particular, for $0\leq j\leq n$ we use
\[
 \mathcal N_{j,\alpha}=\Omega_{V_\alpha}^{[j]},\qquad
 N_j=(N_{j,k})_{k\geq1}.
\]
The notation $\Omega_{\widehat{\mathcal B}}^{[j]}$ will mean this
completed ambient system.  It does not refer to an intrinsic reflexive
differential sheaf on the nonreduced $\mathcal S_k$.

The ambient Hom sheaf for two such objects is
\[
 \mathcal H_{\alpha,MN}
    =\mathcal Hom_{V_\alpha}(\mathcal M_\alpha,\mathcal N_\alpha).
\]
It has the natural \'etale transition maps, and we write
\[
 \widehat{\mathcal H}_{MN}
   =\bigl(\mathcal H_{\alpha,MN}/I_\alpha^k
                         \mathcal H_{\alpha,MN}\bigr)_{k\geq1}
\]
for its completed system.  On an adic complete chart this is the Hom
module of the completed ambient modules: completion is flat for
noetherian rings and Hom commutes with flat base change for finite
presentations.  Composition and trace are induced by their ambient
maps.  Put
\[
 N=\bigoplus_{j=0}^nN_j,\qquad
 \cE=\widehat{\mathcal H}_{NN}.
\]
Thus $\cE$ is completed from the ambient endomorphism sheaf before
any quotient is taken.

Second, a bounded-support coefficient is a coherent sheaf $M$ on some
$\mathcal S_k$, regarded at larger orders by pushforward along the
closed immersions.  For such a coefficient, $N_j\otimes M$ means the ordinary
tensor product $N_{j,k}\otimes_{\cO_{\mathcal S_k}}M$ at any sufficient
order.  Its pullback to $X$ is the ordinary pullback constructed in
Proposition~\ref{str:products}.  This use of coefficients is separate
from the completed ambient systems.  In particular the full truncated
endomorphism sheaf
\[
 \mathcal{E}nd_{\mathcal S_k}\left(\bigoplus_jN_{j,k}\right)
\]
is allowed to have maps that do not lift to ambient endomorphisms.
The natural map from $\cE/I^k\cE$ to this full sheaf may be neither
injective nor surjective.

For ambient objects $M,N$, define their barred Hom sheaf on $\cS$ by
the residual trace pairing:
\begin{equation}\label{str:barred-hom}
 \overline{\Hom}(M,N)=
 \frac{\widehat{\mathcal H}_{MN}|_{\cS}}
 {\ker\left[
   \widehat{\mathcal H}_{MN}|_{\cS}\longrightarrow
       (\widehat{\mathcal H}_{NM}|_{\cS})^\vee,\quad
   a\longmapsto\bigl(b\longmapsto\tr(ba)\bigr)\right]}.
\end{equation}
We put $\overline{\cE}=\overline{\Hom}(N,N)$.  For the fixed natural
tensor expressions and structural maps used here,
Lemma~\ref{str:uniform} below shows that these are locally free and
that their fibers are the trace quotients of the corresponding local
tensor objects.  A residual direct summand is then handled by the
image of its idempotent on these bundles; it is locally a direct
summand bundle as well.
Cyclicity and partial trace make the kernels in
\eqref{str:barred-hom} a tensor ideal, so composition and the indicated
tensor operations descend to these quotients.

Coarse and chart ideals give cofinal filtrations: on a finite chart
the radical of the extended coarse ideal is $I_\alpha$, so suitable
powers contain one another by noetherianity.  A finite cover makes
these comparisons uniform.

\begin{lemma}[Uniform transverse calculations]\label{str:uniform}
On the finite family of product refinements in
Proposition~\ref{str:products}, the ambient sheaves
$\mathcal N_j=\Omega^{[j]}$, the fixed natural mixed reflexive tensor
expressions in them, and their ambient Hom and trace maps have fixed
transverse descriptions.  Their completed systems at order $k$ are
obtained from fixed finite modules over $A_k$ by extension along the
smooth factor.  The same holds for kernels and images of the resulting
fixed transverse maps.  Residual retracts are obtained by applying
their idempotents to the resulting barred bundles.  In particular the barred Hom sheaves in
\eqref{str:barred-hom} are locally free on $\cS$ and their fibers are
the local residual trace Hom spaces.

Consequently a property of these modules and maps that, for each fixed
transverse complete ring, holds at every order $k\geq k_0$ has one
valid order on the whole stratum, provided it is preserved by flat
scalar extension and can be checked by faithfully flat completion.
The completed product descriptions persist after replacing the smooth
parameters by lifts of another coordinate system on the reduced
stratum.
\end{lemma}

\begin{proof}
In our applications the transversal $T_0$ is normal near $t_0$.
Its analytic germ is the transverse factor of a normal product, hence
normal.  Its completion is normal by excellence and is also the
completion of the algebraic local ring; faithful flatness of completion
then descends normality to that ring.  Shrink $T_0$ accordingly.
On a product $\mathbb A^d\times T_0$, restriction to the smooth locus
of the normal transversal gives
\begin{equation}\label{str:cotangent-product}
 \Omega^{[j]}_{\mathbb A^d\times T_0}
 =\bigoplus_{a+b=j}
      \operatorname{pr}_1^*\Omega^a_{\mathbb A^d}
          \otimes\operatorname{pr}_2^*\Omega_{T_0}^{[b]}.
\end{equation}
Both sides are reflexive, so the equality extends across codimension
two.  It remains valid under the common \'etale pullbacks.
Equation~\eqref{str:finite-product} gives the claimed description on
each whole thickened chart.  For finite modules, flat base change
commutes with Hom, as follows from a finite presentation; it also
commutes with kernels and images.  It therefore commutes with double
duals and the fixed reflexive tensor constructions here.  Composition
and contraction are preserved, and trace can be checked on the dense
locally free locus.  The trace pairing on the reduced smooth factor is
thus the scalar extension of a pairing of finite-dimensional complex
vector spaces.  Its kernel and quotient are vector bundles there.
This proves the claim about barred Hom.

We make explicit how the residual tensor category of the full maximal
chart is used in these transverse calculations.  At a product point put
\[
 R=\widehat{\cO}_{T_0,t_0},\qquad
 S=R[[t_1,\ldots,t_d]],\qquad
 E=R^{\oplus d}\oplus\widehat{\Omega_{T_0}^{[1]}}.
\]
The completion of the full chart cotangent module is $E\otimes_RS$.
For finite reflexive tensor objects $P,Q$ generated by $E$, flat
base change gives
\[
 \Hom_S(P\otimes_RS,Q\otimes_RS)
       =\Hom_R(P,Q)\otimes_RS,
\]
and also commutes with their reflexive tensor operations, composition
and trace.  Reduction modulo
$\mathfrak m_S=(\mathfrak m_R,t_1,\ldots,t_d)$ gives the same trace
pairing as reduction over $R$ modulo $\mathfrak m_R$.  Its kernels and
residual Hom spaces therefore agree.  The same flat Hom calculation
compares the analytic local ring of the full chart with its completion.
Idempotent lifting over the henselian rings, as in
Lemmas~\ref{loc:trace-radical} and \ref{loc:split-lift}, identifies
their idempotent-complete residual tensor categories.  Thus the
residual category of the full maximal chart is transported by scalar
extension to the category computed from the transverse module $E$.
Theorem~\ref{loc:slsp} and Corollary~\ref{loc:exterior-constituent}
are applied on the full chart and carried to the transverse
calculation by this identification.

One may now choose an exponent on each of the finitely many
transversals and take their maximum.  Equalities of coherent maps,
surjectivity, and specified powers annihilating a kernel descend from
faithfully flat completions.  This proves precisely the asserted
uniformity for the fixed finite diagrams.

For the last assertion, the stratum ideal in $S$ is
$\mathfrak m_RS$.  If $t'_1,\ldots,t'_d$ restrict to another coordinate
system on $S/\mathfrak m_RS$, substitution $t_i\mapsto t'_i$, fixing
$R$, is an automorphism of this complete ring by the formal inverse
function theorem.  It preserves the stratum ideal.  Completed
differentials and their reflexive tensor maps are functorial for this
automorphism, so it transports the completed Hom, trace and ideal
data as well.
\end{proof}

\begin{proposition}[Compatible direct images on normalized charts]
\label{str:direct-images}
Let $U\to B$ be an integral separated \'etale algebraic neighborhood,
let $\pi:C\to U$ be an integral finite normal Galois quasi-\'etale
chart, and let $h:Z\to C$ be the normalized source chart of
Proposition~\ref{loc:maximal-charts}.  For the global degreewise
isomorphisms $\phi_j:R^jf_*\cO_X\to\Omega_B^{[j]}$ fixed after
Proposition~\ref{prop:background}, there are unique isomorphisms
\[
 \phi_{j,C}:R^jh_*\cO_Z\xrightarrow{\ \sim\ }\Omega_C^{[j]},
 \qquad 0\leq j\leq n,
\]
whose restrictions to $C^\circ=C\times_BB_{\mathrm{reg}}$ are the
ordinary flat pullbacks of $\phi_j$.  The degree-zero map is the
unit.  These isomorphisms are compatible with ordinary \'etale
refinement and with every normalized main overlap whose projections
are \'etale.  In particular they give compatible ambient
identifications $\mathcal N_{j,\alpha}=R^jh_{\alpha*}\cO_{X_\alpha}$
on the charts of Proposition~\ref{str:products}.
\end{proposition}

\begin{proof}
The restriction $C^\circ\to B_{\mathrm{reg}}$ is \'etale.  Indeed it
factors through the smooth $U\times_BB_{\mathrm{reg}}$, and purity
makes the finite quasi-\'etale factor \'etale there.  Its complement
in $C$ has codimension at least two.  Over $C^\circ$ the normalized
source is the ordinary base change $X\times_BC^\circ$: that base change
is smooth and normal, all its components dominate $C^\circ$ by
equidimensionality, and it has the integral generic fiber used in
Proposition~\ref{loc:maximal-charts}.  Proper flat base change and
\'etale pullback of differentials therefore pull $\phi_j$ to
\[
 (R^jh_*\cO_Z)|_{C^\circ}\xrightarrow{\ \sim\ }\Omega_{C^\circ}^j.
\]
This open includes the entire inverse image of the discriminant over
the regular part of $B$; only the singular locus of $B$ has been
removed.

The discrepancy argument in Proposition~\ref{loc:maximal-charts} applies
to any such quasi-\'etale $\pi$, so $C$ is klt.  In particular it is
Cohen--Macaulay and $\Q$-Gorenstein.  The same proposition gives a
projective equidimensional map $h$ with connected fibers and smooth
canonical-trivial source.  Ou's local reflexivity theorem
\cite[Theorem~1.3]{Ou2019} therefore makes $R^jh_*\cO_Z$ reflexive
on $C$.  Both sides of the displayed isomorphism are reflexive and
the omitted set has codimension at least two, so it extends uniquely
to $\phi_{j,C}$.  At $j=0$ it is the Stein unit.

Ordinary \'etale base change preserves both sides.  On a normalized
main overlap with \'etale projections, every component meets the
inverse image of $B_{\mathrm{reg}}$, since a main component meets
the dense common unbranched open.  The two pullbacks there are the
same pullback of the fixed global $\phi_j$.  They consequently
agree on the entire normal overlap by reflexivity.  This also proves
compatibility with deck actions and with further \'etale refinements.
\end{proof}

We finish by recording the scope of spreading.  Choose models over a
finitely generated integral subring of $\C$ for the finite diagram in
Proposition~\ref{str:products}, for the ambient coherent isomorphisms
of Proposition~\ref{str:direct-images}, and for the finite coherent
maps required in Lemma~\ref{str:uniform} at the selected orders.
After localization, the object and overlap charts retain normality,
the stated maps retain \'etaleness near the closed loci, the reduced
stratum charts are smooth, the finite source charts are smooth \'etale
torsors, and the coarse stratum is geometrically integral and projective.
Retain also a fiberwise dense common unbranched open and its finite
torsor identifications.  Keep the ideal equalities and groupoid
identities as identities of the models, and invert the finitely many
deck-group orders.

For each of the pair, triple, quadruple, unit and inverse domains used
above, retain the disjointness of its omitted closed complement from
the corresponding closed support.  There are finitely many such
intersections, all with empty generic fiber, so one further localization
keeps them empty.  This ensures that the quotients by every power of
the marked ideal remain in those domains after reduction.  The common
cross-arrow projections also retain their coverage of the corresponding
closed supports.  Their images are open by \'etaleness, and the closed
subsets of those supports omitted by the relevant unions of images have
empty generic fiber.  There are finitely many such subsets, so one
further localization removes them.  The common
product maps are modeled as \'etale maps of marked pairs.  Their
product descriptions at every finite order then follow from formal
\'etaleness, uniformly in the model.

For any additional fixed finite list of coherent presentations,
kernels, images and identities, generic freeness and coherent base
change retain the required calculations after further localization.
In positive characteristic we use these spread sheaves, direct-image
isomorphisms and algebraic identities.  No positive-characteristic
maximal-cover or residual tensor classification is asserted.
